\chapter{Mathematics}

\section{Matrix Notation for Derivatives}



\subsection*{Scalar-valued function}
Let $x \in \mathbb{R}^N$ be a vector.\\
Let $f(x)$ be a real-valued function:
\[
f : \mathbb{R}^N \to \mathbb{R}
\].


The gradient vector of $f$ at point $x$ is:
\[
\nabla f(x) =
\begin{pmatrix}
\frac{\partial f(x)}{\partial x_1} \\
\vdots \\
\frac{\partial f(x)}{\partial x_N}
\end{pmatrix}
\]


\subsection*{Inner Product}
For $x, y \in \mathbb{R}^N$, the inner product is:
\[
x \cdot y = \sum_{n=1}^N x_n y_n
\]

\subsection*{Vector-valued function}
Let $f : \mathbb{R}^N \to \mathbb{R}^M$ be a vector-valued differentiable function:
\[
f(x) =
\begin{pmatrix}
f_1(x_1, \dots, x_N) \\
f_2(x_1, \dots, x_N) \\
\vdots \\
f_M(x_1, \dots, x_N)
\end{pmatrix}
\]

The derivative (Jacobian matrix) is:
\[
Df(x) =
\begin{pmatrix}
\frac{\partial f_1(x)}{\partial x_1} & \frac{\partial f_1(x)}{\partial x_2} & \cdots & \frac{\partial f_1(x)}{\partial x_N} \\
\frac{\partial f_2(x)}{\partial x_1} & \frac{\partial f_2(x)}{\partial x_2} & \cdots & \frac{\partial f_2(x)}{\partial x_N} \\
\vdots & \vdots & \ddots & \vdots \\
\frac{\partial f_M(x)}{\partial x_1} & \frac{\partial f_M(x)}{\partial x_2} & \cdots & \frac{\partial f_M(x)}{\partial x_N}
\end{pmatrix}
\]



\subsection*{Hessian Matrix}

Let $f : \mathbb{R}^N \to \mathbb{R}$.  
The Hessian matrix is defined as:
\[
H = D^2 f(x) = D[\nabla f(x)]
\]

\[
H(x) =
\begin{pmatrix}
\frac{\partial^2 f(x)}{\partial x_1^2} & \frac{\partial^2 f(x)}{\partial x_1 \partial x_2} & \cdots & \frac{\partial^2 f(x)}{\partial x_1 \partial x_N} \\
\frac{\partial^2 f(x)}{\partial x_2 \partial x_1} & \frac{\partial^2 f(x)}{\partial x_2^2} & \cdots & \frac{\partial^2 f(x)}{\partial x_2 \partial x_N} \\
\vdots & \vdots & \ddots & \vdots \\
\frac{\partial^2 f(x)}{\partial x_N \partial x_1} & \frac{\partial^2 f(x)}{\partial x_N \partial x_2} & \cdots & \frac{\partial^2 f(x)}{\partial x_N^2}
\end{pmatrix}
\]

\section{Homogeneous Functions and Euler's Theorem}

A function $f(x_1, \dots, x_N)$ is said to be homogeneous of degree $r$ if for all $t > 0$:
\[
f(t x_1, \dots, t x_N) = t^r f(x_1, \dots, x_N)
\]

\begin{tcolorbox}[breakable, colback=white]

\subsection*{Theorem 1}

If $f(x_1, \dots, x_N)$ is homogeneous of degree $r$, then for each $n = 1, \dots, N$,
\[
\frac{\partial f(x_1, \dots, x_N)}{\partial x_n}
\]
is homogeneous of degree $r - 1$.


\subsection*{Proof}

Define:
\[
f(t x_1, \dots, t x_N) - t^r f(x_1, \dots, x_N) = 0
\]

Differentiate both sides with respect to $x_n$:
\[
t \frac{\partial f(t x_1, \dots, t x_N)}{\partial tx_n}
- t^r \frac{\partial f(x_1, \dots, x_N)}{\partial x_n} = 0
\]

Thus,
\[
\frac{\partial f(t x_1, \dots, t x_N)}{\partial tx_n}
= t^{r-1} \frac{\partial f(x_1, \dots, x_N)}{\partial x_n}
\]


\end{tcolorbox}

\subsection*{Homothetic Functions}

If $f(x_1, \dots, x_N)$ is homogeneous of degree $r$, and $h(\cdot)$ is an increasing function, then:
\[
h(f(x_1, \dots, x_N))
\]
is called a homothetic function.

The slopes of the level sets of a homothetic function are unchanged along rays through the origin.


\begin{tcolorbox}[breakable]
\subsection*{Exercise 1-1}

Assume that $f(x_1, \dots, x_N)$ and $g(x_1, \dots, x_N)$ are homogeneous
of degrees $r$ and $s$, respectively. \\Check whether the function
$h(x_1, \dots, x_N)$ defined in the following ways is homogeneous.
In the cases in which it is, determine its degree of homogeneity:

\begin{enumerate}
    \item[(i)] $h(x_1, \dots, x_N) = f(x_1^p, \dots, x_N^p)$
    \item[(ii)] $h(x_1, \dots, x_N) = [g(x_1, \dots, x_N)]^p$
    \item[(iii)] $h(x_1, \dots, x_N) \\= f(x_1, \dots, x_N) + g(x_1, \dots, x_N)$
    \item[(iv)] $h(x_1, \dots, x_N) \\= f(x_1, \dots, x_N)\, g(x_1, \dots, x_N)$
    \item[(v)] $h(x_1, \dots, x_N) = \dfrac{f(x_1, \dots, x_N)}{g(x_1, \dots, x_N)}, \qquad g \neq 0$
\end{enumerate}

\end{tcolorbox}

\begin{tcolorbox}[breakable, colback=white]
\subsection*{Theorem 2 (Euler)}

Suppose that $f(x_1, \dots, x_N)$ is homogeneous of degree $r$ and differentiable.
Then, at any $(x_1, \dots, x_N)$, we have
\[
\sum_{n=1}^N x_n \frac{\partial f(x_1, \dots, x_N)}{\partial x_n}
= r f(x_1, \dots, x_N).
\]

Or, in matrix notation,
\[
\nabla f(x) \cdot x = r f(x).
\]

\subsection*{Proof}

By definition, we have
\[
f(t x_1, \dots, t x_N) - t^r f(x_1, \dots, x_N) = 0.
\]

Differentiating with respect to $t$ gives
\[
\sum_{n=1}^N \frac{\partial f(t x_1, \dots, t x_N)}{\partial (t x_n)} x_n
- r t^{r-1} f(x_1, \dots, x_N) = 0.
\]

This holds for any $t$. Evaluating at $t=1$ yields
\[
\sum_{n=1}^N x_n \frac{\partial f(x_1, \dots, x_N)}{\partial x_n}
= r f(x_1, \dots, x_N).
\]

\end{tcolorbox}

\section{Concave and Quasiconcave Functions}

The set $A \subseteq \mathbb{R}^N$ is convex if
\[
\alpha x + (1-\alpha)x' \in A
\]
whenever $x, x' \in A$ and $\alpha \in [0,1]$.\\


The function $f : A \to \mathbb{R}$, defined on a convex set $A \subset \mathbb{R}^N$,
is \textbf{concave} if
\[
f(\alpha x' + (1-\alpha)x) \ge \alpha f(x') + (1-\alpha) f(x)
\]
for all $x, x' \in A$ and all $\alpha \in [0,1]$.

If the inequality is strict for all $x' \neq x$ and all $\alpha \in (0,1)$,
then $f$ is \textbf{strictly concave}.

The function is \textbf{convex} or \textbf{strictly convex} if the inequality is reversed.


\begin{tcolorbox}[breakable, colback=white]
\subsection*{Theorem 3}

A continuously differentiable function $f : A \to \mathbb{R}$ is concave
if and only if
\[
f(x + z) \le f(x) + \nabla f(x) \cdot z
\]
for all $x \in A$ and $z \in \mathbb{R}^N$ such that $x + z \in A$.

The function $f$ is strictly concave if and only if the inequality holds strictly
for all $x \in A$ and $z \neq 0$.

\subsection*{Proof}

Assume concavity. Then for $\alpha \in (0,1)$,
\[
\begin{aligned}
f(x+z) - f(x)
&\le \frac{f(x + \alpha (x + z -x)) - f(x)}{\alpha}\\
&=\frac{f(x + \alpha z) - f(x)}{\alpha}.
\end{aligned}
\]


Taking the limit as $\alpha \to 0$,
\[
\frac{f(x + \alpha z) - f(x)}{\alpha} \to \nabla f(x) \cdot z,
\]
hence
\[
f(x+z) \le f(x) + \nabla f(x) \cdot z.
\]\\

Conversely, assume
\[
f(x+z) \le f(x) + \nabla f(x) \cdot z.
\]

Let
\[
\begin{aligned}
\bar{x} 
&= \alpha (x+z) + (1-\alpha)x \\
&= x + \alpha z.
\end{aligned}
\]

Then,
\[
\begin{aligned}
f(x+z) - f(\bar{x}) 
&\le \nabla f(\bar{x}) \cdot (x+z-\bar{x})\\
&=(1-\alpha)\nabla f(\bar{x}) \cdot z,
\end{aligned}
\]
and
\[
\begin{aligned}
f(x) - f(\bar{x}) 
&\le \nabla f(x) \cdot (x - \bar{x}) \\
&= -\alpha \nabla f(x) \cdot z.
\end{aligned}
\]

Multiply the first inequality by $\alpha$ and the second by $(1-\alpha)$:
\[
\begin{aligned}
&\alpha (f(x+z) - f(\bar{x})) + (1-\alpha)(f(x) - f(\bar{x}))\\
&\le \alpha(1-\alpha)\nabla f(x)\cdot z - (1-\alpha)\alpha \nabla f(x)\cdot z \\
&= 0.
\end{aligned}
\]

Thus,
\[
\alpha f(x+z) + (1-\alpha) f(x) \le f(\bar{x}),
\]
which implies concavity.\\


Suppose that $f$ is strictly concave (and thus concave) and $z \neq 0$. Then,
\[
\begin{aligned}
f(x+z) - f(x)
&< \frac{f(x) - f(\bar{x})}{\alpha} \\
&\le \frac{\nabla f(x)\cdot (\bar{x}-x)}{\alpha}\\
&= \nabla f(x)\cdot z
.
\end{aligned}
\]

If, on the other hand,
\[
f(x+z) < f(x) + \nabla f(x)\cdot z,
\]
then applying the same steps used to prove concavity yields
\[
\alpha f(x+z) + (1-\alpha) f(x) < f(\bar{x}),
\]
and thus $f$ is strictly concave.

\end{tcolorbox}


\vspace{1em}

The $N \times N$ matrix $M$ is \textbf{negative semidefinite} if
\[
z^\top M z \le 0 \quad \text{for all } z \in \mathbb{R}^N.
\]

If the inequality is strict for all $z \neq 0$, then $M$ is
\textbf{negative definite}:
\[
z^\top M z < 0.
\]


\begin{tcolorbox}[breakable, colback=white]
\subsection*{Theorem 4}

A twice continuously differentiable function $f : A \to \mathbb{R}$ is concave
if and only if $D^2 f(x)$ is negative semidefinite for every $x \in A$.

If $D^2 f(x)$ is negative definite for every $x \in A$, then $f$ is strictly concave.

\subsection*{Proof (necessity only)}

Suppose that $f$ is concave. Consider $z$ and $x + \alpha z$, where $\alpha \in \mathbb{R}$.

Define
\[
\phi(\alpha) = f(x + \alpha z).
\]

Taking a Taylor expansion around $\alpha = 0$ gives
\[
\begin{aligned}
&f(x+\alpha z) - f(x) - \nabla f(x)\cdot (\alpha z)\\
&= \frac{\alpha^2}{2} z^\top D^2 f(x + \beta z) z
\end{aligned}
\]
for some $\beta \in [0,\alpha]$. By definition, the left-hand side is nonpositive, hence
\[
\frac{\alpha^2}{2} z^\top D^2 f(x+\beta z) z \le 0.
\]
Taking sufficiently small $\alpha$ and $\beta$ yields
\[
z^\top D^2 f(x) z \le 0. \quad 
\]

\end{tcolorbox}



\textbf{Notice.} Strict concavity does not necessarily imply negative definiteness.  
Example: $f(x) = -x^4$.

\vspace{1em}

We have the expansion:
\[
\begin{aligned}
f(x+\alpha z)
&= f(x) + \nabla f(x)\cdot (\alpha z)
+ \frac{\alpha^2}{2} z^\top D^2 f(x+\beta z) z\\
&= f(x) + \nabla f(x)\cdot (\alpha z)
+ \frac{\alpha^2}{2} z^\top D^2 f(x) z + R\\
&\approx f(x) + \nabla f(x)\cdot (\alpha z)
+ \frac{\alpha^2}{2} z^\top D^2 f(x) z.
\end{aligned}
\]

If $f$ is continuous on $[x, x+\alpha z]$, then by the Mean Value Theorem,
there exists $\beta z$ such that
\[
f(x+\alpha z) = f(x) + f'(x+\beta z)(\alpha z).
\]

Its second-order version (Lagrange form of the remainder) is:
\[
f(x+\alpha z)
= f(x) + f'(x)(\alpha z)
+ \frac{(\alpha z)^2}{2} f''(x+\beta z).
\]


\begin{tcolorbox}[breakable]
\subsection*{Exercise 1-2}

Define the function
\[
f(x_1,x_2) = A\left(\delta x_1^{\rho} + (1-\delta)x_2^{\rho}\right)^{1/\rho},
\]
where $\rho \neq 0$, $0 \le \delta \le 1$, $A>0$, $x_1>0$, $x_2>0$.

This is the CES function. Examine whether it is concave or convex.
\end{tcolorbox}

\begin{tcolorbox}[breakable]

\subsection*{Exercise 1-3}

Prove that the function
\[
f(x) = \left(x_1^2 + \cdots + x_N^2\right)^{1/2}
\]
is convex.

\textit{Hint: This is a norm, and every norm is convex.}

\end{tcolorbox}


\vspace{1em}

The function $f : A \to \mathbb{R}$ defined on a convex set $A \subset \mathbb{R}^N$
is \textbf{quasiconcave} if its upper contour sets
\[
\{x \in A : f(x) \ge t\}
\]
are convex.

That is, if $f(x) \ge t$ and $f(x') \ge t$, then
\[
f(\alpha x + (1-\alpha)x') \ge t
\]
for any $t \in \mathbb{R}$, $x,x' \in A$, and $\alpha \in [0,1]$.

If the inequality is strict whenever $x \neq x'$ and $\alpha \in (0,1)$,
then $f$ is \textbf{strictly quasiconcave}.

The function is \textbf{quasiconvex} if its lower contour sets are convex.


\begin{tcolorbox}[breakable, colback=white]
\subsection*{Theorem 5}

A continuously differentiable function $f : A \to \mathbb{R}$ is
\textbf{quasiconcave} if and only if
\[
\nabla f(x)\cdot (x' - x) \ge 0
\]
\[
\quad \text{whenever } f(x') \ge f(x), \ \forall x,x' \in A.
\]

If
\[
\nabla f(x)\cdot (x' - x) > 0
\]
whenever $f(x') \ge f(x)$ and $x' \ne x$, then $f$ is
\textbf{strictly quasiconcave}.

Conversely, if $f$ is strictly quasiconcave and $\nabla f(x) \ne 0$
for all $x \in A$, then
\[
\nabla f(x)\cdot (x' - x) > 0
\]
\[
\text{whenever } f(x') \ge f(x), \ x' \ne x.
\]

\subsection*{Proof (necessity only)}

Assume that $f$ is quasiconcave. If $f(x') \ge f(x)$ and $\alpha \in [0,1]$, define
\[
g(\alpha) = f(\alpha x' + (1-\alpha)x).
\]

Then,
\[
g(\alpha) \ge \min\{f(x), f(x')\} = f(x) = g(0).
\]

Hence $g(\alpha)$ is increasing at $\alpha=0$, and
\[
g'(0) = \nabla f(x)\cdot (x' - x) \ge 0. \quad \blacksquare
\]



\end{tcolorbox}

\vspace{1em}

The condition
\[
\nabla f(x)\cdot (x' - x) > 0
\]
whenever $f(x') \ge f(x)$ and $x' \ne x$
implies that the angle between the vectors $(x'-x)$ and $\nabla f(x)$ is acute.

\vspace{1em}

\textbf{Remark.} The condition $\nabla f(x) \ne 0$ is necessary.  
For example, $f(x) = x^3$ is strictly quasiconcave, but
\[
\nabla f(0) = 0,
\]
so
\[
\nabla f(x)\cdot (x' - x) = 0 \quad \text{when } x=0.
\]

\vspace{1em}

If $f$ is concave, then it is quasiconcave:
\[
f(\alpha x + (1-\alpha)x')
\ge \alpha f(x) + (1-\alpha)f(x')
\ge \min\{f(x), f(x')\}.
\]

Thus, if $f(x) \ge t$ and $f(x') \ge t$, then
\[
f(\alpha x + (1-\alpha)x') \ge t.
\]

\begin{tcolorbox}[breakable]

\subsection*{Exercise 1-4}

Assume $x \in \mathbb{R}$. Prove that any monotonic function of $x$
is both quasiconvex and quasiconcave.


\end{tcolorbox}


\begin{tcolorbox}[breakable, colback=white]

\subsection*{Theorem 6}

Let $f : A \to \mathbb{R}$ be a twice continuously differentiable function,
and assume that $\nabla f(x) \ne 0$ for every $x \in A$.

Then $f$ is \textbf{quasiconcave} if and only if $D^2 f(x)$ is negative semidefinite
on the subspace
\[
\{ z \in \mathbb{R}^N : \nabla f(x)\cdot z = 0 \}
\]
for every $x \in A$.

That is,
\[
z^\top D^2 f(x) z \le 0
\quad \text{whenever } \nabla f(x)\cdot z = 0, \ \forall x \in A.
\]

If $D^2 f(x)$ is negative definite on the subspace
\[
\{ z \in \mathbb{R}^N : \nabla f(x)\cdot z = 0,\ z \ne 0 \}
\]
for every $x \in A$, then $f$ is \textbf{strictly quasiconcave}.

\end{tcolorbox}


\vspace{1em}

\textbf{Remark.} If the condition $\nabla f(x) \ne 0$ for every $x$ is not imposed,
then
\[
z^\top D^2 f(x) z \le 0 \quad \text{whenever } \nabla f(x)\cdot z = 0
\]
is necessary but not sufficient for quasiconcavity.

Example:
\[
f(x) = (x_1 - 1)^2 (x_2 - 1)^2.
\]

\vspace{1em}

\subsection*{Summary 1 (Concavity)}

\begin{itemize}
    
    \item $f(x)$ is concave \\
    $\iff z^\top D^2 f(x) z \le 0 \quad \forall z \in \mathbb{R}^N$
    \item $f(x)$ is strictly concave \\
    $\impliedby z^\top D^2 f(x) z < 0 \quad \forall z \in \mathbb{R}^N, z \ne 0$ 
\end{itemize}

\subsection*{Summary 2 (Quasiconcavity)}

\begin{itemize}
    \item $f(x)$ is quasiconcave \\
    $\iff z^\top D^2 f(x) z \le 0$\\
    whenever $\nabla f(x)\cdot z = 0$ and $\nabla f(x) \ne 0$
    \item $f(x)$ is strictly quasiconcave \\
    $\impliedby z^\top D^2 f(x) z < 0$\\
    whenever $\nabla f(x)\cdot z = 0$, $z \ne 0$
\end{itemize}



Let $M$ be a $T \times S$ matrix. 
Then ${}_t M_{s}$ is the $t \times s$ submatrix of $M$ where only the first $t \le T$ rows and $s \le S$ columns are retained.

Let $M$ be an $N \times N$ symmetric matrix. Then $M$ is \textbf{negative definite}
if and only if
\[
(-1)^r \det({}_r M_r) > 0
\qquad \text{for every } r=1,\dots,N,
\]
where $M_r$ denotes the $r \times r$ leading principal minor of $M$.

\vspace{1em}

\textbf{Example 4.} For $f(x_1,x_2)$,
\[
f_{11} < 0, \qquad
\begin{vmatrix}
f_{11} & f_{12} \\
f_{21} & f_{22}
\end{vmatrix} > 0.
\]
Equivalently,
\[
f_{11} < 0, \qquad f_{11}f_{22} - (f_{12})^2 > 0.
\]

\vspace{1em}

$M$ is \textbf{negative semidefinite} if and only if
\[
(-1)^r \det({}_r M_r^\pi) \ge 0
\]
for every $r=1,\dots,N$ and for every permutation $\pi$ of the indices
$\{1,\dots,N\}$.



\vspace{1em}

Let $B$ be an $N \times S$ matrix with $S \le N$ and rank equal to $S$.
Then $M$ is negative definite in the space
\[
\{z \in \mathbb{R}^N : B \cdot z = 0,\ z \ne 0\}
\]
that is,
\[
z^\top M z < 0
\quad \text{for all } z \in \mathbb{R}^N \text{ with } B\cdot z = 0,\ z\ne 0,
\]
if and only if
\[
(-1)^r
\det
\begin{pmatrix}
0 & B_r \\
B_r^\top & {}_r M_r
\end{pmatrix}
> 0
\qquad \text{for } r=S+1,\dots,N.
\]

Moreover, $M$ is negative semidefinite in the space
\[
\{z \in \mathbb{R}^N : B\cdot z = 0\}
\]
if and only if
\[
(-1)^r
\det
\begin{pmatrix}
0 & B_r^\pi \\
(B_r^\pi)^\top & {}_r M_r^\pi
\end{pmatrix}
\ge 0
\]
for $r=S+1,\dots,N$ and every permutation $\pi$, where $B_r^\pi$
is the matrix formed by permuting only the rows of $B_r$ according to $\pi$.

\vspace{1em}

\textbf{Example 5.} Since $\nabla f(x)\ne 0$, assume that at least $f_1 \ne 0$.
Then $f(x)$ is strictly quasiconcave if $D^2 f(x_1,x_2)$ is negative definite
in the subspace
\[
\{z \in \mathbb{R}^2 : \nabla f(x)\cdot z = 0,\ z\ne 0\}
\qquad \forall x=(x_1,x_2).
\]
Here $S=1$ and $r=2$.

By the theorem above, this holds if and only if
\[
\begin{vmatrix}
0 & f_1 & f_2 \\
f_1 & f_{11} & f_{12} \\
f_2 & f_{21} & f_{22}
\end{vmatrix}
> 0.
\]

Under the constraint
\[
f_1 z_1 + f_2 z_2 = 0,
\]
we have
\[
[z_1 \ z_2]
\begin{bmatrix}
f_{11} & f_{12} \\
f_{21} & f_{22}
\end{bmatrix}
\begin{bmatrix}
z_1 \\ z_2
\end{bmatrix}
=
\left(-\frac{z_2 f_2}{f_1},\ z_2\right)
\begin{bmatrix}
f_{11} & f_{12} \\
f_{21} & f_{22}
\end{bmatrix}
\begin{bmatrix}
-\frac{z_2 f_2}{f_1} \\ z_2
\end{bmatrix}
\]
\[
=
\left(\frac{f_2}{f_1}\right)^2 f_{11} z_2^2
-2\left(\frac{f_2}{f_1}\right) f_{12} z_2^2
+ f_{22} z_2^2,
\]
which is negative when $D^2 f(x_1,x_2)$ is negative definite in the subspace
\[
\{z \in \mathbb{R}^2 : \nabla f(x)\cdot z = 0,\ z\ne 0\}.
\]

\begin{tcolorbox}[breakable]

\subsection*{Exercise 1-5}

Prove that the Cobb--Douglas function
\[
f(x) = A x_1^{\alpha_1} x_2^{\alpha_2} \cdots x_N^{\alpha_N}
\]
is

\begin{enumerate}
    \item[(i)] always quasiconcave,
    \item[(ii)] strictly concave if $\alpha = \alpha_1+\cdots+\alpha_N < 1$,
    \item[(iii)] concave if and only if $\alpha \le 1$.
\end{enumerate}

\textit{Hint: Taking a log transformation may help.}

\end{tcolorbox}


\section{Implicit Function Theorem}

A system of $N$ equations depending on $N$ endogenous variables
\[
x = (x_1,\dots,x_N)
\]
and $M$ parameters
\[
q = (q_1,\dots,q_M)
\]
is given by
\[
f_n(x_1,\dots,x_N;\, q_1,\dots,q_M)=0,
\qquad n=1,\dots,N.
\]



Suppose that
\[
\bar{x}=(\bar{x}_1,\dots,\bar{x}_N)\in A \subset \mathbb{R}^N
\quad \text{and} \quad
\bar q=(\bar q_1,\dots,\bar q_M)\in B \subset \mathbb{R}^M
\]
satisfy the above system of equations.

Is it possible to solve for
\[
{x}=({x}_1,\dots,{x}_N)
\]
as a function of
\[
{q}=({q}_1,\dots,{q}_M)
\]
locally around $\bar x$ and $\bar{q}$?\\

An open neighborhood (or open ball) of a point $x\in\mathbb{R}^N$ is
\[
A'=\{x'\in\mathbb{R}^N:\|x'-x\|<\varepsilon\}
\]
for some scalar $\varepsilon>0$, where
\[
\|x'-x\|
=
\sqrt{(x_1'-x_1)^2+\cdots+(x_N'-x_N)^2}.
\]


\begin{tcolorbox}[breakable, colback=white]
\subsection*{Theorem 7 (Implicit Function Theorem)}

Suppose that every equation $f_n(\cdot)$ is continuously differentiable
with respect to its $N+M$ variables, and that
\[
f_n(\bar{x};\bar q)=0
\qquad \forall n.
\]

If
\[
\det
\begin{pmatrix}
\frac{\partial f_1(\bar{x};\bar q)}{\partial x_1} & \cdots & \frac{\partial f_1(\bar{x};\bar q)}{\partial x_N} \\
\vdots & \ddots & \vdots \\
\frac{\partial f_N(\bar{x};\bar q)}{\partial x_1} & \cdots & \frac{\partial f_N(\bar{x};\bar q)}{\partial x_N}
\end{pmatrix}
\neq 0,
\]
then the system can be locally solved at $(\bar{x},\bar q)$ by implicitly defined
functions
\[
\eta_n : B' \to A'
\]
that are continuously differentiable.

Moreover,
\[
D_q \eta(\bar q)
=
-\bigl[D_x f(\bar{x};\bar q)\bigr]^{-1} D_q f(\bar{x};\bar q).
\]

\subsection*{Proof}

Proving the existence of $\eta_n:B'\to A'$ is very technical.
However, note that the condition in the theorem is a full-rank condition.
By this condition, we can adjust $x$ in any direction to restore the equilibrium
when $q$ changes. Therefore, a change in $q$ can be completely reflected by a
change in $x$, and there exists $\eta_n:B'\to A'$.

Since
\[
f(\eta(q);q)=0
\qquad \forall q\in B',
\]
we differentiate to obtain
\[
D_x f(\bar x;\bar q)\, D_q \eta(\bar q) + D_q f(\bar x;\bar q)=0.
\]

By the full-rank condition, $D_x f(\bar x;\bar q)$ is invertible. Therefore,
\[
D_q \eta(\bar q)
=
-\bigl[D_x f(\bar x;\bar q)\bigr]^{-1} D_q f(\bar x;\bar q).
\qquad \blacksquare
\]

\end{tcolorbox}

\textbf{Example 7.}  
If $M=N=1$, then
\[
\frac{d\eta(\bar q)}{dq}
=
-\frac{\partial f(\bar x;\bar q)/\partial q}{\partial f(\bar x;\bar q)/\partial x}.
\]

\begin{tcolorbox}[breakable]

\subsection*{Exercise 1-6}

Consider the system
\[
u^3 + xu - 2y - z = 0
\]
and
\[
u - zv^3 + xy + xv = 0.
\]

Can you define $u$ and $v$ as differentiable functions of $x,y,z$
in a neighborhood of
\[
(x,y,z,u,v)=(0,1,-1,1,-1)?
\]

If the answer is yes, then compute the partial derivatives of
$u$ and $v$ with respect to $x$, $y$, and $z$.

\end{tcolorbox}


\section{Continuous Functions and Compact Sets}

A sequence in $\mathbb{R}^N$ assigns to every positive integer $m=1,2,\dots$
a vector $x_m \in \mathbb{R}^N$, denoted by $\{x_m\}$.

\medskip

The sequence $\{x_m\}$ converges to $x \in \mathbb{R}^N$ if for every $\varepsilon>0$,
there exists an integer $M_\varepsilon$ such that
\[
\|x_m - x\| < \varepsilon
\quad \text{whenever } m > M_\varepsilon.
\]
We write
\[
\lim_{m\to\infty} x_m = x
\quad \text{or} \quad x_m \to x.
\]

\medskip

Consider a domain $X \subset \mathbb{R}^N$. A function
\[
f : X \to \mathbb{R}
\]
is \textbf{continuous} if for all $x \in X$ and every sequence $x_m \to x$,
we have
\[
f(x_m) \to f(x).
\]

Equivalently, $f$ is continuous if for every $\varepsilon>0$, there exists
$\delta(\varepsilon)>0$ such that
\[
\|x - x_0\| < \delta(\varepsilon)
\quad \Rightarrow \quad
|f(x) - f(x_0)| < \varepsilon.
\]
A small change in $x$ results in a small change in $f(x)$


\medskip

Fix a set $X \subset \mathbb{R}^N$.

A set $A \subset X$ is \textbf{open} (relative to $X$) if for every $x \in A$,
there exists $\varepsilon>0$ such that
\[
\|x' - x\| < \varepsilon, \ x' \in X \ \Rightarrow \ x' \in A.
\]

That is, for every point in $A$, there exists an open ball around it
contained in $A$. Every point in A remains in A under small perturbations.

\medskip

A set $A \subset X$ is \textbf{closed} (relative to $X$) if its complement
$X \setminus A$ is open.

Equivalently, $A$ is closed if for every sequence $x_m \to x \in X$,
with $x_m \in A$ for all $m$, we have
\[
x \in A.
\]

\medskip

\textbf{Properties:}
\begin{itemize}
    \item The union of any number of open sets is open.
    \item The intersection of a finite number of open sets is open.
    \item The intersection of any number of closed sets is closed.
    \item The union of a finite number of closed sets is closed.
\end{itemize}

\vspace{10pt}

A set $A \subset \mathbb{R}^N$ is \textbf{bounded} if there exists $r \in \mathbb{R}$
such that
\[
\|x\| < r \quad \text{for every } x \in A.
\]

The set $A \subset \mathbb{R}^N$ is \textbf{compact} if it is bounded and closed.

\subsection*{Theorem 8 (Weierstrass)}

Suppose that $f : X \to \mathbb{R}$ is continuous and $X \subset \mathbb{R}^N$
is nonempty and compact. Then $f$ attains a maximum and a minimum on $X$.

\vspace{1em}

The Weierstrass theorem follows from the following lemmata.

\subsection*{Lemma 1}

If $f : X \to \mathbb{R}$ is continuous on $X \subset \mathbb{R}^N$, and $X$ is compact,
then $f(X)$ is also compact.

\subsection*{Lemma 2}

If $A \subset \mathbb{R}$ is compact, then
\[
\sup A \in A \quad \text{and} \quad \inf A \in A,
\]
so the maximum and minimum of $A$ are well defined.

\section{Convex Sets and Separating Hyperplanes}

Given a set $B \subset \mathbb{R}^N$, the \textbf{convex hull} of $B$,
denoted by $\operatorname{co}(B)$, is the smallest convex set containing $B$.

That is,
\[
\operatorname{co}(B)
=
\bigcap \{ C \subset \mathbb{R}^N : C \text{ is convex and } B \subset C \}.
\]


\subsection*{(continued)}

$\operatorname{co}(B)$ is the set of all possible convex combinations of elements of $B$:
\[
\operatorname{co}(B)
=
\left\{
\sum_{j=1}^J a_j x_j
\;\middle|\;
x_j \in B,\ a_j \ge 0,\ \sum_{j=1}^J a_j = 1
\right\}.
\]

\vspace{1em}

Given $p \in \mathbb{R}^N$ with $p \ne 0$, and $c \in \mathbb{R}$,
the \textbf{hyperplane} generated by $p$ and $c$ is
\[
H_{p,c} = \{ z \in \mathbb{R}^N : p \cdot z = c \}.
\]

The sets
\[
\{ z \in \mathbb{R}^N : p \cdot z \ge c \}
\quad \text{and} \quad
\{ z \in \mathbb{R}^N : p \cdot z \le c \}
\]
are called, respectively, the half-space above and below the hyperplane $H_{p,c}$.

\subsection*{Theorem 9 (Separating Hyperplane)}

\begin{itemize}
    \item[(i)] Suppose that $B \subset \mathbb{R}^N$ is convex and closed,
    and that $x \notin B$. Then there exists $p \in \mathbb{R}^N$, $p \ne 0$,
    and $c \in \mathbb{R}$ such that
    \[
    p \cdot x > c
    \quad \text{and} \quad
    p \cdot y < c \quad \forall y \in B.
    \]

    \item[(ii)] Suppose that the convex sets $A,B \subset \mathbb{R}^N$ are disjoint.
    Then there exists $p \in \mathbb{R}^N$, $p \ne 0$, and $c \in \mathbb{R}$ such that
    \[
    p \cdot x \ge c \quad \forall x \in A,
    \qquad
    p \cdot y \le c \quad \forall y \in B.
    \]
\end{itemize}

That is, there exists a hyperplane that separates $A$ and $B$
(Minkowski separating hyperplane theorem).

\subsection*{Theorem 10 (Supporting Hyperplane)}

Let $B \subset \mathbb{R}^N$ be a convex set and let $x$ be a boundary point of $B$.
Then there exists at least one supporting hyperplane at $x$.

That is, there exists $p \in \mathbb{R}^N$, $p \ne 0$, such that
\[
p \cdot x \ge p \cdot y
\quad \forall y \in B.
\]




\section{Unconstrained Maximization}

A vector $x \in \mathbb{R}^N$ is a \textbf{local maximizer} of $f(\cdot)$
if there exists an open neighborhood $A \subset \mathbb{R}^N$ of $x$ such that
\[
f(x) \ge f(x') \quad \text{for every } x' \in A.
\]

If
\[
f(x) \ge f(x') \quad \text{for every } x' \in \mathbb{R}^N,
\]
then $x$ is a \textbf{global maximizer}.

\medskip

Suppose that $f(x)$ is differentiable and that $x \in \mathbb{R}^N$
is a local maximizer or local minimizer of $f(\cdot)$. Then
\[
\frac{\partial f(x)}{\partial x_n} = 0 \quad \forall n,
\qquad \text{or} \qquad
\nabla f(x) = 0.
\]

Otherwise, $f$ would be locally increasing or decreasing
in the direction of some coordinate axis.

\subsection*{Theorem 11}

Suppose that $f : \mathbb{R}^N \to \mathbb{R}$ is twice continuously differentiable
and that
\[
\nabla f(x) = 0.
\]

\begin{itemize}
    \item[(i)] If $x$ is a local maximizer, then $D^2 f(x)$ is negative semidefinite.
    \item[(ii)] If $D^2 f(x)$ is negative definite, then $x$ is a local maximizer.
\end{itemize}

\subsection*{Proof}

\textbf{(i)} Let $z \in \mathbb{R}^N$ be any direction and $\varepsilon > 0$.
Consider
\[
\phi(\varepsilon) = f(x + \varepsilon z).
\]

A Taylor expansion around $\varepsilon=0$ gives
\[
f(x+\varepsilon z) - f(x)
=
\varepsilon \nabla f(x)\cdot z
+ \frac{1}{2}\varepsilon^2 z^\top D^2 f(x) z + R.
\]

Since $\nabla f(x)=0$,
\[
f(x+\varepsilon z) - f(x)
=
\frac{1}{2}\varepsilon^2 z^\top D^2 f(x) z + R.
\]

If $x$ is a local maximizer, then
\[
f(x+\varepsilon z) - f(x) \le 0.
\]

Taking the limit as $\varepsilon \to 0$, we obtain
\[
z^\top D^2 f(x) z \le 0.
\]

\medskip

\textbf{(ii)} If
\[
z^\top D^2 f(x) z < 0 \quad \forall z \ne 0,
\]
then for sufficiently small $\varepsilon > 0$,
\[
f(x+\varepsilon z) < f(x),
\]
so $x$ is a local maximizer. \quad $\blacksquare$




\subsection*{Example 8}

Consider the function
\[
f(x) = x^3.
\]

Then
\[
\frac{df(x)}{dx}\bigg|_{x=0} = 0,
\]
but $x=0$ is neither a maximum nor a minimum.

Thus, a negative semidefinite $D^2 f(x)$ does not guarantee a maximum.

\begin{tcolorbox}

\subsection*{Exercise 7}

Consider the function
\[
f(x,y) = x^2 - xy + y^2 + 3x - 2y + 1.
\]

Find the minimum of $f$ and prove that it is the global minimum.

\end{tcolorbox}


\section{Equality Constrained Maximization}

Consider the problem ($N \ge M$):
\[
\max_{x \in \mathbb{R}^N} f(x)
\]
subject to
\[
g_1(x) = b_1,\quad \dots,\quad g_M(x) = b_M.
\]

The set of all feasible points is called the \textbf{constraint set}:
\[
C = \{x \in \mathbb{R}^N : g_m(x) = b_m,\ m=1,\dots,M\}.
\]

\subsection*{Theorem 12 (The Theorem of Lagrange)}

Suppose that the objective function $f$ and the constraint functions
$g_1,\dots,g_M$ are differentiable, and that $x \in C$ is a local maximizer.

If the gradients $\nabla g_1(x), \dots, \nabla g_M(x)$ are linearly independent,
then there exist scalars $\lambda_1,\dots,\lambda_M$ such that
\[
\nabla f(x)
=
\sum_{m=1}^M \lambda_m \nabla g_m(x).
\]

Equivalently,
\[
\nabla f(x) = \lambda_1 \nabla g_1(x) + \cdots + \lambda_M \nabla g_M(x).
\]

\subsection*{Lagrangian}

Define the Lagrangian function:
\[
\mathcal{L}(x,\lambda)
=
f(x)
+ \sum_{m=1}^M \lambda_m \bigl(b_m - g_m(x)\bigr).
\]

The first-order conditions are:
\[
\nabla_x \mathcal{L}(x,\lambda) = 0,
\qquad
\frac{\partial \mathcal{L}}{\partial \lambda_m} = 0 \quad (m=1,\dots,M).
\]

That is,
\[
\nabla f(x)
=
\sum_{m=1}^M \lambda_m \nabla g_m(x),
\quad
g_m(x)=b_m.
\]


\subsection*{(continued)}

Suppose that $x \in C$ is a local constrained maximizer.
Assume also that the $M \times N$ matrix
\[
\begin{pmatrix}
\frac{\partial g_1(x)}{\partial x_1} & \cdots & \frac{\partial g_1(x)}{\partial x_N} \\
\vdots & \ddots & \vdots \\
\frac{\partial g_M(x)}{\partial x_1} & \cdots & \frac{\partial g_M(x)}{\partial x_N}
\end{pmatrix}
\]
has rank $M$.

Then there exist scalars $\lambda_m \in \mathbb{R}$ (Lagrange multipliers),
one for each constraint, such that
\[
\frac{\partial f(x)}{\partial x_n}
=
\sum_{m=1}^M \lambda_m \frac{\partial g_m(x)}{\partial x_n},
\qquad n=1,\dots,N.
\]

Equivalently,
\[
\nabla f(x)
=
\sum_{m=1}^M \lambda_m \nabla g_m(x).
\]

\vspace{1em}

If a pair $(x,\lambda)$ satisfies
\[
g_m(x)=b_m \quad \text{and} \quad
\nabla f(x)=\sum_{m=1}^M \lambda_m \nabla g_m(x),
\]
we say that $(x,\lambda)$ satisfies the \textbf{first-order necessary conditions}
(FOC) for equality-constrained optimization.

\medskip

The rank condition is called the \textbf{constraint qualification condition},
and it implies that the constraints are locally independent.

\vspace{1em}

\textbf{Interpretation.}  
At the optimum, the gradient of $f$ lies in the span of the gradients of the constraints.
That is, there is no feasible direction along which $f$ can increase.

\vspace{1em}

\subsection*{Example 9 (Elimination Approach)}

Consider
\[
\max f(x_1,x_2)
\quad \text{s.t.} \quad g(x_1,x_2)=b.
\]

The rank condition implies that at least one of $g_1,g_2$ is nonzero.
Assume $g_1 \ne 0$.

Then there exists $\lambda$ such that
\[
f_1(x)=\lambda g_1(x),
\qquad
\lambda=\frac{f_1(x)}{g_1(x)}.
\]

Define the constraint:
\[
\phi(x_1,x_2)=g(x_1,x_2)-b=0.
\]

Since $g_1 \ne 0$, we can locally solve for
\[
x_1 = h(x_2).
\]

Thus the problem becomes an unconstrained problem:
\[
\max f(h(x_2), x_2).
\]

The first-order condition is
\[
f_1(x) h'(x_2) + f_2(x)=0.
\]

From the implicit function theorem,
\[
h'(x_2) = -\frac{g_2(x)}{g_1(x)}.
\]

Substituting,
\[
f_1(x)\left(-\frac{g_2(x)}{g_1(x)}\right) + f_2(x)=0
\quad \Rightarrow \quad
\frac{f_1(x)}{g_1(x)} = \frac{f_2(x)}{g_2(x)}.
\]

Thus
\[
\nabla f(x) = \lambda \nabla g(x).
\]

\vspace{1em}

\begin{tcolorbox}
\subsection*{Exercise 8}

Let
\[
f(x_1,x_2) = -x_1^2 - x_2^2,
\qquad
g(x_1,x_2)=x_2^2 - x_1^2=0.
\]

Show that $(0,0)$ is the unique maximizer.
Check the constraint qualification condition at $(0,0)$.

\end{tcolorbox}


\vspace{1em}

\subsection*{Example 10}

If $N=2$, $M=1$, and $b=0$, then the condition
\[
\nabla f(x)=\lambda \nabla g(x)
\]
implies that the vectors $\nabla f(x)$ and $\nabla g(x)$ lie on the same line.

That is, at the optimum, the gradient of the objective function is parallel
to the gradient of the constraint.


\subsection*{(continued)}

The point $x$ is a constrained maximum under the constraint $g(x)=0$.
This is because the curve $g(x)=0$ is tangent at $x$ to the level curve
$f(x)=f(x)$.

The vector $\nabla g(x)$ is perpendicular to the constraint curve $g(x)=0$,
and the vector $\nabla f(x)$ is perpendicular to the level curve $f(x)=f(x)$.

\medskip

\textbf{Why is the gradient perpendicular to the level curve?}

Consider a level curve
\[
f(x_1,x_2)=k.
\]

Treat $x_2$ as a function of $x_1$ and differentiate implicitly:
\[
f_1(x_1,x_2) + f_2(x_1,x_2)\frac{dx_2}{dx_1}=0.
\]

Thus, the slope of the tangent line is
\[
\frac{dx_2}{dx_1}
=
-\frac{f_1(x_1,x_2)}{f_2(x_1,x_2)}.
\]

On the other hand, the slope of the gradient direction is
\[
\frac{\Delta x_2}{\Delta x_1}
=
\frac{f_2(x_1,x_2)}{f_1(x_1,x_2)}.
\]

Since these are negative reciprocals, the gradient $\nabla f(x)$ is
perpendicular to the tangent line of the level curve.

\medskip

\textbf{General case.}

For $N \ge 3$ and $M \ge 2$, the theorem requires that
\[
\nabla f(x)
\in
\operatorname{span}\{\nabla g_1(x), \dots, \nabla g_M(x)\}.
\]

That is,
\[
\nabla f(x)
=
\sum_{m=1}^M \lambda_m \nabla g_m(x).
\]

\medskip

Define the Lagrangian:
\[
\mathcal{L}(x,\lambda)
=
f(x) + \sum_{m=1}^M \lambda_m (b_m - g_m(x)).
\]

The Theorem of Lagrange implies that if $x$ is a local constrained maximizer
(and the constraint qualification holds), then there exist multipliers
$\lambda_1,\dots,\lambda_M$ such that all partial derivatives of the Lagrangian
are zero:

\[
\frac{\partial \mathcal{L}(x,\lambda)}{\partial x_n} = 0
\quad \text{for } n=1,\dots,N,
\]

\[
\frac{\partial \mathcal{L}(x,\lambda)}{\partial \lambda_m} = 0
\quad \text{for } m=1,\dots,M.
\]

\medskip

Equivalently, the first-order conditions are:
\[
\nabla f(x)
=
\sum_{m=1}^M \lambda_m \nabla g_m(x),
\qquad
g_m(x)=b_m.
\]


\subsection*{(continued)}

Do the same thing for the contour of $g(x)$:
\[
\frac{dx_2}{dx_1} = -\frac{g_1(x_1,x_2)}{g_2(x_1,x_2)}.
\]

Then
\[
\frac{d^2 x_2}{dx_1^2}
=
\frac{d}{dx_1}\left(-\frac{g_1}{g_2}\right)
=
\frac{
g_2^2 g_{11} - 2 g_1 g_2 g_{12} + g_1^2 g_{22}
}{g_2^3}.
\]

The condition of relative convexity requires
\[
\frac{d^2 x_2}{dx_1^2}\bigg|_{f}
>
\frac{d^2 x_2}{dx_1^2}\bigg|_{g}.
\]

That is,
\[
\frac{
f_2^2 f_{11} - 2 f_1 f_2 f_{12} + f_1^2 f_{22}
}{f_2^3}
-
\frac{
g_2^2 g_{11} - 2 g_1 g_2 g_{12} + g_1^2 g_{22}
}{g_2^3}
> 0.
\]

By substituting the first-order conditions
\[
f_1 = \lambda g_1,
\qquad
f_2 = \lambda g_2,
\]
into the above condition, we obtain
\[
g_2^2(\lambda g_{11} - f_{11})
-2g_1g_2(\lambda g_{12} - f_{12})
+g_1^2(\lambda g_{22} - f_{22}) > 0.
\]

It can be shown that this condition is equivalent to
\[
\begin{vmatrix}
L_{11} & L_{12} & -g_1 \\
L_{21} & L_{22} & -g_2 \\
-g_1 & -g_2 & 0
\end{vmatrix}
> 0,
\]
where
\[
L_{11} = f_{11} - \lambda g_{11}, \qquad
L_{12} = f_{12} - \lambda g_{12},
\]
\[
L_{21} = f_{21} - \lambda g_{21}, \qquad
L_{22} = f_{22} - \lambda g_{22}.
\]

\begin{tcolorbox}

    \subsection*{Exercise 9}

Consider the problem
\[
\max \ (\min)\ x^2 + y^2 + z^2
\]
subject to
\[
x + 2y + z = 1
\]
and
\[
2x - y - 3z = 4.
\]

Obtain the solution using the Lagrange method.
Did you find a maximum, a minimum, or neither?

\end{tcolorbox}

\section{Inequality Constrained Maximization}

Consider the problem
\[
\max_{x \in \mathbb{R}^N} f(x)
\]
subject to
\[
g_m(x)=b_m, \qquad m=1,\dots,M,
\]
\[
h_k(x)\le c_k, \qquad k=1,\dots,K.
\]

Again, let $C$ denote the constraint set.

We say that the \textbf{constraint qualification} is satisfied at $x \in C$
if the constraints that hold with equality at $x$ are linearly independent.

\subsection*{Theorem 14 (Kuhn--Tucker)}

Suppose that $x \in C$ is a local maximizer and that the constraint qualification is satisfied.

Then there exist multipliers
\[
\lambda_m \in \mathbb{R}, \qquad m=1,\dots,M,
\]
for the equality constraints, and
\[
\mu_k \in \mathbb{R}_+, \qquad k=1,\dots,K,
\]
for the inequality constraints, such that
\[
\frac{\partial f(x)}{\partial x_n}
=
\sum_{m=1}^M \lambda_m \frac{\partial g_m(x)}{\partial x_n}
+
\sum_{k=1}^K \mu_k \frac{\partial h_k(x)}{\partial x_n},
\qquad n=1,\dots,N,
\]
and
\[
\mu_k\bigl(c_k - h_k(x)\bigr)=0,
\qquad k=1,\dots,K.
\]

\medskip

Define the Lagrangian function
\[
\mathcal{L}(x,\lambda,\mu)
=
f(x)
+
\sum_{m=1}^M \lambda_m \bigl(b_m - g_m(x)\bigr)
+
\sum_{k=1}^K \mu_k \bigl(c_k - h_k(x)\bigr).
\]


\subsection*{(continued)}

The Lagrangian function is defined as
\[
\mathcal{L}(x,\lambda,\mu)
=
f(x)
+
\sum_{m=1}^M \lambda_m \bigl(b_m-g_m(x)\bigr)
+
\sum_{k=1}^K \mu_k \bigl(c_k-h_k(x)\bigr).
\]

The Kuhn--Tucker theorem says that if $x$ is a local constrained maximizer
(and the constraint qualification is satisfied), then there exist values
$\lambda_m$ and $\mu_k$ such that
\[
\frac{\partial \mathcal{L}(x,\lambda,\mu)}{\partial x_n}
=
\frac{\partial f(x)}{\partial x_n}
-
\sum_{m=1}^M \lambda_m \frac{\partial g_m(x)}{\partial x_n}
-
\sum_{k=1}^K \mu_k \frac{\partial h_k(x)}{\partial x_n}
=0
\]
for every $n=1,\dots,N$,

\[
\frac{\partial \mathcal{L}(x,\lambda,\mu)}{\partial \lambda_m}
=
b_m-g_m(x)=0
\qquad \text{for every } m=1,\dots,M,
\]
and
\[
\mu_k\bigl(c_k-h_k(x)\bigr)=0
\qquad \text{for every } k=1,\dots,K.
\]

\medskip

The \textbf{complementary slackness condition}
\[
\mu_k\bigl(c_k-h_k(x)\bigr)=0
\]
implies that $\mu_k=0$ for any constraint $k$ that does not bind, that is,
for any constraint such that
\[
h_k(x) < c_k.
\]

\medskip

A formal proof of the Kuhn--Tucker theorem can be derived from the
Theorem of Lagrange. Consider the problem
\[
\mathcal{L}(x,\mu)
=
f(x)+\sum_{k=1}^K \mu_k\bigl(c_k-h_k(x)\bigr).
\]

Then for some values of $\mu_k$,
\[
\frac{\partial \mathcal{L}}{\partial x_n}
=
\frac{\partial f(x)}{\partial x_n}
-
\sum_{k=1}^K \mu_k \frac{\partial h_k(x)}{\partial x_n}
=0
\qquad \text{for every } n=1,\dots,N,
\]
and
\[
\mu_k\bigl(c_k-h_k(x)\bigr)=0
\qquad \text{for every } k=1,\dots,K.
\]

Take auxiliary variables $y_k$ $(k=1,\dots,K)$, and define
\[
\Phi(x,y,\mu)
=
f(x)+\sum_{k=1}^K \mu_k\bigl(c_k-h_k(x)-y_k\bigr).
\]

This turns the inequality-constrained problem into an equality-constrained one
with
\[
c_k-h_k(x)-y_k=0.
\]

Then the Theorem of Lagrange implies
\[
\frac{\partial \Phi}{\partial x_n}
=
\frac{\partial f(x)}{\partial x_n}
-
\sum_{k=1}^K \mu_k \frac{\partial h_k(x)}{\partial x_n}
=0
\qquad \text{for every } n=1,\dots,N,
\]
\[
\frac{\partial \Phi}{\partial y_k}=-\mu_k=0,
\]
and
\[
\frac{\partial \Phi}{\partial \mu_k}
=
c_k-h_k(x)-y_k=0
\qquad \text{for every } k=1,\dots,K.
\]

If $\mu_k=0$, then the complementary slackness condition is satisfied directly.
If $\mu_k>0$, then necessarily
\[
y_k = c_k-h_k(x)=0,
\]
so again the Kuhn--Tucker conditions are satisfied.

Thus the Lagrange conditions for the auxiliary equality-constrained problem
lead to the Kuhn--Tucker conditions.

\subsection*{A Simple Case}

Consider the case $K=1$, $M=0$, and $c=0$. Then the constraint is
\[
h(x)\le 0.
\]

For example, if $h(x)=-x$, then the problem is equivalent to
\[
x \ge 0.
\]

The Lagrangian is
\[
\mathcal{L}(x,\mu)=f(x)+\mu x.
\]

The Kuhn--Tucker conditions are
\[
f'(x)+\mu=0,
\qquad
\mu x=0,
\qquad
\mu \ge 0,
\qquad
x\ge 0.
\]

There are three possible cases:

\begin{itemize}
    \item[(a)] \textbf{Interior solution:} $x>0$.  
    Then $\mu=0$, and so
    \[
    f'(x)=0.
    \]

    \item[(b)] \textbf{Corner solution with binding constraint:} $x=0$ and $\mu>0$.  
    Then
    \[
    f'(x)<0.
    \]

    \item[(c)] \textbf{Corner solution with zero multiplier:} $x=0$ and $\mu=0$.  
    Then
    \[
    f'(x)=0.
    \]
\end{itemize}

\subsection*{Example 12}

Consider
\[
\max U(x,y)= y + a\ln x,
\qquad a>0.
\]

The Lagrangian is
\[
\mathcal{L}(x,y,\mu,\mu_x,\mu_y)
=
y + a\ln x
+
\mu( I - p_x x - p_y y )
+
\mu_x x
+
\mu_y y.
\]

The first-order conditions are
\[
\mathcal{L}_x
=
\frac{a}{x}-\mu p_x+\mu_x=0,
\]
\[
\mathcal{L}_y
=
1-\mu p_y+\mu_y=0.
\]




\subsection*{(continued)}

The first-order conditions are:
\[
(1)\quad \mathcal{L}_x = \frac{a}{x} - \mu p_x + \mu_x = 0,
\]
\[
(2)\quad \mathcal{L}_y = 1 - \mu p_y + \mu_y = 0,
\]
\[
(3)\quad \mu \bigl(I - p_x x - p_y y\bigr) = 0,
\]
\[
(4)\quad \mu_x x = 0,
\]
\[
(5)\quad \mu_y y = 0.
\]

\medskip

Suppose that $I - p_x x - p_y y > 0$ and $\mu = 0$.
Then from (1),
\[
\frac{a}{x} + \mu_x = 0,
\]
which implies $\mu_x < 0$, a contradiction since $\mu_x \ge 0$.
Thus, the budget constraint must be binding:
\[
p_x x + p_y y = I.
\]

\medskip

We consider the following four cases:

\begin{itemize}
    \item[(i)] $x=0,\ y=0$.  
    This is not feasible since the budget constraint is not binding.

    \item[(ii)] $x=0,\ y=\frac{I}{p_y}$.  
    From (2),
    \[
    1 - \mu p_y = 0 \quad \Rightarrow \quad \mu = \frac{1}{p_y}.
    \]
    Substituting into (1),
    \[
    \frac{a}{x} - \mu p_x + \mu_x = 0,
    \]
    which implies $\frac{a}{x} = \infty$, hence not feasible.

    \item[(iii)] $x>0,\ y=0$.  
    Then $x=\frac{I}{p_x}$ and $\mu_x=0$.

    From (1),
    \[
    \frac{a}{x} - \mu p_x = 0
    \quad \Rightarrow \quad
    \mu = \frac{a}{I}.
    \]

    From (2),
    \[
    1 - \mu p_y + \mu_y = 0
    \quad \Rightarrow \quad
    1 \le \frac{a p_y}{I}.
    \]

    \item[(iv)] $x>0,\ y>0$.  
    Then $\mu_x = \mu_y = 0$.

    From (1) and (2),
    \[
    \frac{a}{x} = \mu p_x, \qquad 1 = \mu p_y.
    \]

    Thus,
    \[
    \mu = \frac{1}{p_y},
    \quad
    x = \frac{a p_y}{p_x}.
    \]

    Using the budget constraint,
    \[
    y = \frac{I}{p_y} - a.
    \]

    This is feasible if
    \[
    I > a p_y.
    \]
\end{itemize}

\medskip

\textbf{Conclusion.}

There are two possible solutions:

\begin{itemize}
    \item[(i)] Corner solution:
    \[
    x = \frac{I}{p_x}, \quad y = 0
    \quad \text{if } I \le a p_y.
    \]

    \item[(ii)] Interior solution:
    \[
    x = \frac{a p_y}{p_x}, \quad y = \frac{I}{p_y} - a
    \quad \text{if } I > a p_y.
    \]
\end{itemize}

\vspace{1em}

A formal statement of second-order sufficient conditions for inequality
constraints is given in the appendix.

The following theorem provides an example where the SOCs are satisfied.



\subsection*{Theorem 15}

Suppose that the constraint set $C$ of the inequality-constrained problem
is convex and that the objective function $f(\cdot)$ is strictly quasiconcave.
Then there is a unique global constrained maximizer.

\subsection*{Proof}

If $x$ and $x' \ne x$ were both global constrained maximizers, then for any
$\alpha \in (0,1)$, the point
\[
x''=\alpha x+(1-\alpha)x'
\]
would be feasible, that is, $x'' \in C$, since $C$ is convex.

By the strict quasiconcavity of $f(\cdot)$,
\[
f(x'') = f(\alpha x+(1-\alpha)x') > f(x)=f(x'),
\]
which contradicts the assumption that both $x$ and $x'$ are global constrained maximizers.
\hfill $\blacksquare$

\begin{tcolorbox}

\subsection*{Exercise 10}

Find a solution to the problem
\[
\max f(x,y,z)=x+y+z
\]
subject to
\[
g_1(x,y,z)=x^2+y^2+z^2 \le 1,
\]
\[
g_2(x,y,z)=x-y-z \le 1,
\]
\[
g_3(x,y,z)=-x \le 0.
\]

\end{tcolorbox}

\begin{tcolorbox}

\subsection*{Exercise 11}

Consider the following cost-minimization problem, where
\[
x_1 + x_2^{\delta} \ge y
\]
represents the production technology:
\[
\min \; w_1 x_1 + w_2 x_2
\]
subject to
\[
x_1 \ge 0,\qquad x_2 \ge 0,\qquad x_1 + x_2^{\delta} - y \ge 0.
\]

We may transform the problem into a maximization problem:
\[
\max \; -w_1 x_1 - w_2 x_2
\]
subject to
\[
x_1 \ge 0,\qquad x_2 \ge 0,\qquad x_1 + x_2^{\delta} - y \ge 0.
\]

\begin{enumerate}
    \item[(i)] Set up the Lagrangian and derive the first-order conditions
    of the Kuhn--Tucker theorem.
    \item[(ii)] Derive the explicit solution(s).
\end{enumerate}

\end{tcolorbox}

\section{Envelope Theorem}

We often encounter optimization problems such as
\[
\max_x f(x,a)
\qquad \text{s.t.} \qquad g(x,a)=0,
\]
where $x$ is a vector of choice variables and $a$ is a vector of parameters
that may enter the objective function, the constraints, or both.

Clearly, the solution to the problem depends on $a$; we may denote it by
\[
x(a).
\]

We can define a new function $V(a)$, which gives the value achieved by the
objective function when $x$ is chosen optimally subject to the constraints.
This is called the \textbf{maximum value function}, and it is defined by
\[
V(a)=f(x(a),a)=\max_x f(x,a)\quad \text{s.t.} \quad g(x,a)=0.
\]

Many times we want to know how the maximum value of a function changes as
one or more parameters of the problem change. The envelope theorem can be
used to analyze the behavior of the maximum value function.

\subsection*{Theorem 16 (Envelope)}

Consider a maximization problem
\[
\max_x f(x,a)
\qquad \text{s.t.} \qquad g(x,a)=0,
\]
and suppose that the objective function and constraint are continuously
differentiable in $a$.

Let $x(a)$ solve the problem, and suppose that it is also continuously
differentiable in $a$.

Let
\[
\mathcal{L}(x,a,\lambda)=f(x,a)+\lambda g(x,a)
\]
be the Lagrangian, and let $x(a)$ and $\lambda(a)$ solve the problem.
Finally, let $V(a)$ be the associated maximum value function.

Then the envelope theorem states that
\[
\frac{dV(a)}{da}
=
\frac{\partial \mathcal{L}}{\partial a}\Big|_{x(a),\,a,\,\lambda(a)}.
\]


\subsection*{Proof of Theorem 16}

Evaluate the partial derivative of the Lagrangian with respect to $a_j$
at the point $(x(a),\lambda(a))$:
\[
\frac{\partial \mathcal{L}}{\partial a_j}
=
\frac{\partial f(x(a),a)}{\partial a_j}
+
\lambda(a)\frac{\partial g(x(a),a)}{\partial a_j}.
\]

Using the chain rule, we obtain
\[
\frac{\partial V(a)}{\partial a_j}
=
\sum_{i=1}^N
\frac{\partial f(x(a),a)}{\partial x_i}
\frac{\partial x_i(a)}{\partial a_j}
+
\frac{\partial f(x(a),a)}{\partial a_j}.
\]

From the first-order conditions,
\[
\frac{\partial f(x(a),a)}{\partial x_i}
+
\lambda(a)\frac{\partial g(x(a),a)}{\partial x_i}
=0,
\quad i=1,\dots,N.
\]

Substituting into the summation term,
\[
\frac{\partial V(a)}{\partial a_j}
=
\lambda(a)
\sum_{i=1}^N
\frac{\partial g(x(a),a)}{\partial x_i}
\frac{\partial x_i(a)}{\partial a_j}
+
\frac{\partial f(x(a),a)}{\partial a_j}.
\]

From the constraint $g(x(a),a)=0$, we obtain
\[
\sum_{i=1}^N
\frac{\partial g(x(a),a)}{\partial x_i}
\frac{\partial x_i(a)}{\partial a_j}
+
\frac{\partial g(x(a),a)}{\partial a_j}
=0.
\]

Thus,
\[
\frac{\partial V(a)}{\partial a_j}
=
\lambda(a)\frac{\partial g(x(a),a)}{\partial a_j}
+
\frac{\partial f(x(a),a)}{\partial a_j}
=
\frac{\partial \mathcal{L}}{\partial a_j}.
\quad \blacksquare
\]

\subsection*{Example 13}

Consider
\[
\max_{x_1,x_2} x_1 x_2
\quad \text{s.t.} \quad a - 2x_1 - 4x_2 = 0.
\]

The Lagrangian is
\[
\mathcal{L}
=
x_1 x_2 + \lambda (a - 2x_1 - 4x_2).
\]

FOCs:
\[
x_2 - 2\lambda = 0,
\qquad
x_1 - 4\lambda = 0,
\qquad
a - 2x_1 - 4x_2 = 0.
\]

Solution:
\[
x_1(a) = \frac{a}{4},
\qquad
x_2(a) = \frac{a}{8},
\qquad
\lambda(a) = \frac{a}{16}.
\]

Value function:
\[
V(a)
=
x_1(a)x_2(a)
=
\frac{a}{4}\cdot \frac{a}{8}
=
\frac{a^2}{32}.
\]

Thus,
\[
\frac{dV(a)}{da} = \frac{a}{16}.
\]

Using the envelope theorem:
\[
\frac{dV(a)}{da} = \lambda(a),
\]
which matches the result.

\subsection*{Interpretation of $\lambda$}

Consider
\[
\mathcal{L}(x,a,b,\lambda)=f(x,a)+\lambda\bigl(b-g(x,a)\bigr),
\]
where $b$ measures the strength of the constraint.

Then
\[
\frac{\partial V(a,b)}{\partial b}
=
\lambda(a,b).
\]

Thus, the Lagrange multiplier represents the marginal impact of relaxing
the constraint on the optimal value.

\subsection*{Unconstrained Case}

For
\[
\max_x f(x,a),
\]
we have
\[
\frac{dV(a)}{da}
=
\frac{\partial f(x,a)}{\partial a}\Big|_{x=x(a)}.
\]

\subsection*{Example 14 (First-Price Auction)}

Let $F_v(v)$ be the cumulative distribution of values,
and $N$ the number of bidders.

Let $b(v)$ be the bidding function.
The probability of winning is
\[
G(b(v)) = F_v(v)^{N-1}.
\]

A bidder's payoff is
\[
\pi(v,b(v)) = G(b(v))\bigl(v - b(v)\bigr).
\]

By the envelope theorem,
\[
\frac{d\pi(v)}{dv}
=
G(b(v)) = F_v(v)^{N-1}.
\]

Integrating over $v$ yields the equilibrium bidding function.




\subsection*{(continued)}

Using the value $\pi(0)=0$ as the constant of integration, we obtain
\[
\pi(v,b(v))
=
\pi(0)
+
\int_0^v F_v(x)^{N-1} dx
=
\int_0^v F_v(x)^{N-1} dx,
\]
since a bidder with $v=0$ never wins, so $\pi(0,b(0))=0$.

\medskip

Combining
\[
\pi(v,b(v)) = G(b(v))(v-b(v)),
\quad
\frac{d\pi(v)}{dv} = G(b(v)),
\]
and the integral expression above, we derive the optimal bidding function:
\[
b(v)
=
v - \frac{\int_0^v F_v(x)^{N-1} dx}{F_v(v)^{N-1}}.
\]

Thus,
\[
b(v) < v,
\]
i.e., bidders shade their bids below their true valuation.

\medskip

\textbf{Uniform distribution case.}

Suppose
\[
F_v(v) = \frac{v}{\bar v},
\quad \text{for } v \in (0,\bar v).
\]

Then
\[
\int_0^v F_v(x)^{N-1} dx
=
\int_0^v \left(\frac{x}{\bar v}\right)^{N-1} dx
=
\frac{1}{\bar v^{N-1}} \int_0^v x^{N-1} dx
=
\frac{1}{\bar v^{N-1}} \cdot \frac{v^N}{N}.
\]

Also,
\[
F_v(v)^{N-1}
=
\left(\frac{v}{\bar v}\right)^{N-1}
=
\frac{v^{N-1}}{\bar v^{N-1}}.
\]

Substituting into the bidding function,
\[
b(v)
=
v - \frac{\frac{v^N}{N\bar v^{N-1}}}{\frac{v^{N-1}}{\bar v^{N-1}}}
=
v - \frac{v}{N}
=
\frac{N-1}{N}v.
\]